\scanpage{19}
Soit $X$ un \unsure{schém.} algébrique lisse sur \deleted{\unsure{$k$}} un corps $k$ isomorphe à $\mathbf C$ [de façon non précisée]. Alors
\diagram{\begin{tikzcd}[ampersand replacement=\&]
 A^p(X) \arrow[r,"P_{\mathrm{alg}}"] \& \underline{\underline H}^{2p}(X,\Omega_X^\bullet)
\end{tikzcd}}
est bien défini. Le choix d'un isom $\sigma$ $k\xleftarrow{\sim}\mathbf C$ fixé permet de définir
\margin{\unsure{cano} = \uline{De Rham} + \uline{\unsure{Gr.}}}
\diagram{\begin{tikzcd}[ampersand replacement=\&]
 P_\sigma^{(p)}:\quad A^p(X) \arrow[r,"{\unsure{\mathrm{cano}}\,\unsure{P_{\mathrm{top}}^{an}}\;\deleted{\unsure{P_\sigma^{an}}}}"] \& \underline{\underline H}^{2p}(X,\Omega_X^\bullet)
\end{tikzcd}}
qui se factorise par un sous-réseau $E_\sigma^{\unsure{2i|2p}}=H^{2p}(X_\sigma,\mathbf Z)$ de $\underline{\underline H}^{2p}(X,\Omega_X^\bullet)$. On a
\[
 P_\sigma^{(p)}=\left(\frac{1}{\sigma(2i\pi)}\right)^p P_{\mathrm{alg}},
\]
ce qui montre que $P_\sigma^{(p)}$ dépend beaucoup de $\sigma$, \inserted{$\pi$ étant \uline{transcendant}}. Dans le cas où
\diagram{\begin{tikzcd}[ampersand replacement=\&]
 A^p(X) \arrow[r,"P_\sigma^{(p)}","\unsure{\sim}"'] \& H^{2p}(X_\sigma,\mathbf Z)
\end{tikzcd}}
est un isom (ex: $p=1$, $X=\mathbf P^1$) on voit de $\unsure{=}$ que le réseau $E_\sigma^{\unsure{2i|2p}}$, égal à $\displaystyle\frac{1}{\sigma(2i\pi)}P_{\mathrm{alg}}(A^p(X))$, dépend beaucoup de $\sigma$. \uline{NB} \unsure{Remplaçant} $\sigma$ par $\bar\sigma$ change $P_\sigma^{(p)}$ en $-P_\sigma^{(p)}$, donc ne change pas les $E_\sigma^i=H^{\unsure{2i|i}}(X_\sigma,\mathbf Z)$

\mspage{18}
